\section{Combining Setups, Risk Logic, and Execution Reality}

\subsection{Combination logic: one setup can become another}

A major theme across the lessons is that setups are not isolated boxes. A bounce attempt can
turn into a breakout if the defending participant disappears. A breakout can turn into a false
breakout if follow-through fails. An iceberg can first act as absorption, then later become the
fuel for a strong move once exhausted. A spring trade can be thought of as a bounce where the
stretch makes otherwise modest resistance meaningful.

This implies that future detectors should not be framed as independent yes/no classifiers only.
A state-machine interpretation is more faithful:

\begin{itemize}[leftmargin=1.4em]
  \item identify the current structural state,
  \item watch the event that tests the state,
  \item observe which branch the market takes next,
  \item re-label the active hypothesis as new information arrives.
\end{itemize}

\subsection{A practical confidence model}

The materials support a simple human-style confidence stack:

\begin{enumerate}[leftmargin=1.4em]
  \item \term{Location}: is the event happening where others care?
  \item \term{Participant evidence}: is there a visible or hidden actor worth respecting?
  \item \term{Tape confirmation}: are aggressive traders succeeding or failing?
  \item \term{Context}: is the day supportive of this setup family?
  \item \term{Exit clarity}: do we know where we are wrong?
\end{enumerate}

If several of these are weak, the example is ``unclean'' even if the screenshot looks similar to
an ideal training example.

\subsection{Per-trade risk}

The risk book is conceptually conservative in the right way. Losses are normal, risk is finite,
and position size must be defined against a maximum acceptable loss, not against wishful profit.
In the transcript logic, bounce trades are attractive precisely because they often permit larger
size with a smaller, more local stop. Breakout trades are more dangerous precisely because local
support is less reliable and slippage can be severe.

Practical rules implied by the combined materials:

\begin{itemize}[leftmargin=1.4em]
  \item Every trade should have an invalidation point before entry.
  \item If invalidation is not local and legible, size should shrink sharply or the trade should be
        skipped.
  \item If the trade thesis fails behaviorally, exit quickly rather than wait for a prettier stop.
\end{itemize}

\subsection{Per-day risk and stopping rules}

The risk book strongly supports daily-loss limits. The core idea is simple: once the planned
daily drawdown is reached, the session ends. This matters more in discretionary intraday trading
than in slower styles because emotional state degrades quickly after repeated short-term losses.

The transcripts support this indirectly: markets change fast, what worked yesterday may not work
today, and adaptation matters. A daily stop is therefore not only emotional protection. It is
also a recognition that the trader may be operating with the wrong model for the day.

\subsection{Why overtrading happens}

The materials imply several causes:

\begin{itemize}[leftmargin=1.4em]
  \item seeing one familiar visual cue and ignoring missing confirming reasons,
  \item revenge after a failed breakout or missed move,
  \item trying to trade limit-state or news chaos without a real edge,
  \item treating activity as productivity.
\end{itemize}

Automation can help here, not by removing discretion entirely, but by enforcing structural
filters: daily stop switches, maximum concurrent risk, trade-count limits, and ``no trade unless
required evidence is present'' rules.

\subsection{Execution trade-offs: market versus limit}

The knowledge base is implicitly very execution-aware. In fast setups, especially breakouts,
marketable entry may be the only realistic way to participate. But market entry increases
slippage risk and can destroy reward-to-risk if the move is already extended. Limit entry
improves price but risks missing the move or never getting filled.

For a research platform, this means backtests cannot assume magical fills:

\begin{itemize}[leftmargin=1.4em]
  \item queue position matters,
  \item fill probability matters,
  \item spread state matters,
  \item the difference between touching a price and actually trading there matters.
\end{itemize}

\subsection{Low-liquidity traps}

The old Russian data in the repository contains many sparse books. The transcripts also discuss
instruments whose books are effectively empty for trading purposes. In such environments,
displayed liquidity can be misleading, stop quality is poor, and slippage dominates. A serious
research engine must classify low-liquidity conditions explicitly rather than letting all
setups operate as if queue quality were stable.

\subsection{Differences between crypto and Russian market execution}

The source materials are rooted more strongly in the Russian prop / DOM environment, but the
repository explicitly targets both Russian markets and crypto. The following translation is
therefore necessary:

\begin{itemize}[leftmargin=1.4em]
  \item Crypto often offers richer continuous trading, more accessible public trade feeds, and more
        frequent regime shifts across the full day.
  \item Russian exchange trading is more session-structured, more sensitive to local schedule,
        auction phases, and instrument-specific exchange rules.
  \item Limit-state logic, settlement variants, and guide-instrument relationships can be more
        structurally important in the Russian market than in typical crypto pairs.
\end{itemize}

The framework should share one event model across both domains while keeping venue-specific
features and session metadata explicit.
